\section*{Proyecto de Vida --- Sharon Liz Quino}

\subsection*{Quién soy}
Mi nombre es Sharon Liz Quino, tengo 17 años y estoy en 6to de secundaria, en el BTH de Sistemas Informáticos. Estoy terminando una etapa muy linda de mi vida para empezar otra más grande en la universidad.

Soy una estudiante con futuro y con mucha ambición. Quiero superarme, aprender cosas nuevas y tener un buen futuro para mí y para mi familia. Me gusta la tecnología, programar y hacer soluciones con la computadora. Por eso quiero seguir la carrera de Sistemas. Hoy entiendo que no importa la carrera que uno tome, ahora todo se trata de tecnología, y lo que aprendo ahora me va a servir toda la vida.

\subsection*{Cómo soy y qué quiero mejorar}
Tengo metas claras y ganas de salir adelante, pero sé que tengo cosas que mejorar. A veces me cuesta organizarme y me distraigo fácil. Quiero aprender de mis errores y avanzar poco a poco. Este proyecto me enseñó a ser más responsable: aprendí a revisar bien los datos, a probar muchas veces y a no dejar las cosas a medias. Cuando algo fallaba, buscaba el error con paciencia. Eso me ayudó a no darme por vencida y a cumplir lo que prometo.

\subsection*{Mis metas a corto, mediano y largo plazo}
A corto plazo quiero terminar el colegio con buenas notas, titularme del BTH y entrar a la universidad en Sistemas. También quiero que la página del Módulo siga funcionando bien para que nadie haga filas. A mediano plazo quiero aprender más programación y cuidado de datos. A largo plazo quiero ser ingeniera, tener un buen trabajo en tecnología y apoyar siempre a la sociedad.

\subsection*{Lo que este proyecto significa para mí}
Este proyecto lo hicimos para apoyar a nuestro Módulo. Vimos que las familias sufrían en las filas y quisimos ayudar con tecnología. Ahora la inscripción es más fácil, más rápida y más justa. Eso me hace sentir orgullosa. Este proyecto me ayuda a crecer: aprendí a guardar los datos en orden y a probar la página con muchas personas a la vez. Preguntar a 60 postulantes y hablar con 5 funcionarios me ayudó a perder el miedo a preguntar. Trabajar con Heidy me enseñó a trabajar en equipo: ella hizo las pantallas y yo cuidé los datos. Recuperar la contraseña en menos de 2 minutos y repartir los paralelos por colegio es experiencia real para mi futura carrera de Sistemas.
